% P10-HITL, 15.09.2026: Entscheidungszweck und belegte Zustandswirkung präzisiert.
% P10-4b, 14.09.2026: Begriffs-/Beleganschlüsse und abgegrenzte Prüfaussagen präzisiert.
% M08, 11.09.2026: Entwurfsbegründung konzentriert; Mechanismus in 5.6.
% E46-01 und die Unterscheidung von Gestaltungsziel/Realisierung erhalten.
\section{Persistenter Projektzustand und kontrollierte Mutation}
\label{sec:projektzustand-kontrollierte-mutation}

\emph{Erweitertes Nutzungsziel.} Die in
Abschnitt~\ref{sec:iteratives-vorgehen-evidenz} beschriebene Erweiterung
führte über die einmalige Artefakterzeugung hinaus: Weitere
Projektkommunikation sollte mit dem bereits Erarbeiteten in Beziehung
gesetzt werden. Bei einem zweiten Meeting ist etwa zu klären, ob eine
Aussage neu ist, ein bestehendes Requirement verändert oder ihm
widerspricht. Dafür muss nachvollziehbar sein, was aktuell gilt und
worauf dieser Stand beruht. Diese Fragen ergeben sich aus dem
angestrebten Kreislauf; sie werden hier nicht aus einem gescheiterten
Lauf abgeleitet. Die Form der Zustandshaltung war damit noch offen.

\emph{Entwurfsentscheidung.} Die frühen Dateien in einzelnen
Laufverzeichnissen besaßen weder laufübergreifende fachliche
Identitäten noch explizite Beziehungen oder eine gemeinsame
Änderungshistorie. Für die fortlaufende Bearbeitung wurde deshalb ein
fachlicher Projektzustand mit diesen Eigenschaften eingeführt.
Er wird außerhalb einzelner Läufe geführt: Ein gespeicherter
Workflowzustand ermöglicht die Wiederaufnahme der Ausführung,
beantwortet aber nicht, welche fachlichen Inhalte aktuell gelten
(Abschnitt~\ref{sec:zustand-persistenz-projektwissen}). Die Trennung
von Identität und veränderlichem Inhalt knüpft an die Prinzipien des
Requirements Managements in
Abschnitt~\ref{subsec:requirements-management-evolution} an.
Dokumente und operative Systeme stellen den Bestand dar; seine
maßgebliche Repräsentation liegt im fachlichen Zustandsmodell.
Daraus ergibt sich die fünfte Designanforderung:

\textbf{A5:} Der autorisierte fachliche Projektzustand muss unabhängig
von einzelnen Chatverläufen, Workflowläufen und Dokumenten persistent
geführt werden.

\emph{Abgeleitete Konsequenz.} Mit dem gemeinsamen Bestand entsteht
eine weitere Gestaltungsaufgabe: Transkripte, Klärungsantworten und
externe Rückmeldungen können Änderungen auslösen. Auch ein fachlich
plausibler Vorschlag legt noch nicht fest, ob eine bisherige Anforderung
weiter gelten oder durch eine neue ersetzt werden soll. Die menschliche
Autorisierung soll diese verbindliche Festlegung ermöglichen und zugleich
eine Prüfung der vorgeschlagenen Änderung erlauben. Getrennte
Schreibwege würden die Durchsetzung von Autorisierungs- und
Integritätsregeln verteilen und könnten auseinanderlaufen.
Das Gestaltungsziel verbindet deshalb die menschliche Autorisierung
aus Abschnitt~\ref{sec:governance-autorisierung-kontrollen} mit einer
gemeinsamen Speicheroperation, der \emph{Speichernaht}. Vorschläge
sollen pfadspezifisch geprüft und autorisiert werden; die
Speicheroperation prüft vor der Persistenz strukturelle Invarianten,
etwa zulässige Endpunkte einer Beziehung. Diese Prüfung ersetzt
keine fachliche Entscheidung. Den Ablauf und die Einschränkung des
initialen Schreibpfads beschreibt Abschnitt~\ref{sec:governance-mutation}.

Dies ist eine abgeleitete Realisierungsanforderung, keine
nachträgliche Begründung für den persistenten Zustand. Ein späterer
Entwicklungsfall verdeutlichte das Risiko fehlerhafter Schreibregeln:
Eine deterministische Zuordnung übernahm einen unbestätigten
Herkunftshinweis als Beziehung und brachte dadurch fehlerhafte
Relationen in den Bestand. Nach ihrer Erkennung durch eine
Integritätsmessung wurde der Schreibweg entfernt, die Ableitung auf
bereits bestätigte Zustandsinformationen begrenzt und der Bestand
kontrolliert bereinigt (E46-01). Determinismus allein garantiert
somit keine fachlich richtige Regel. Daraus ergibt sich die sechste
Designanforderung:

\textbf{A6:} Kein vorgesehener Eingang darf den maßgeblichen
Projektzustand außerhalb eines gemeinsamen autorisierten und
invariantengeschützten Mutationspfads verändern.

\emph{Einordnung und Grenzen.} Zustandshaltung und Kontrollen
erzeugen Schema-, Pflege- und Interaktionsaufwand. Der geführte
Bestand ist zudem keine objektive Wahrheit: Autorisierte Inhalte
können falsch sein, strukturell zulässige Beziehungen fachlich
unzutreffend. Kapitel~\ref{chap:loesungsarchitektur} konkretisiert den
Zustand als \emph{Core}; Kapitel~\ref{chap:maf-realisierung} zeigt die
Realisierung. Abschnitt~\ref{sec:eval-kontrollen} untersucht ausgewählte
Kontrollfälle. Noch zu klären ist, wie neue Informationen relativ
zum bestehenden Bestand eingeordnet werden; dies behandelt der
folgende Abschnitt.
